\section{All-order Hessian-sector certificates}\label{app:sectors}

This appendix proves Theorem~\ref{thm:sectors}.  Throughout, $N=n-1$ and
\[
 \pi_k=\frac{\binom{N-1}{k-1}}{2^{N-1}},\qquad 1\le k\le N, \tag{A.1}\label{A.1}
\]
is the complete active-rank law.  Its rank chain has
\[
 R_{k,k+1}=\frac{N-k}{2N},\qquad
 R_{k,k-1}=\frac{k-1}{2N}.                                  \tag{A.2}\label{A.2}
\]
Let $h$ solve $(I-R)h=k^{-1}-c_0$ modulo constants and put
$d_k=h_k-h_{k+1}$ for $1\le k<N$, with $d_0=d_N=0$.
Subtracting adjacent Poisson equations
gives
\[
 (N-k)d_k-(k-1)d_{k-1}
 =2N\left(\frac1k-c_0\right),\qquad 1\le k<N.              \tag{A.3}\label{A.3}
\]

Solving this recurrence gives, for $1\le k<N$,
\[
 \frac{2}{k}-d_k=
 \frac{2^{2-N}\sum_{r=k}^{N-1}\binom{N-1}{r}}
  {(N-1)\binom{N-2}{k-1}},
 \qquad
 \frac{2(N-2)}{Nk}\le d_k\le \frac{2}{k}.                \tag{A.3a}\label{A.3a}
\]
For the lower bound, after clearing the positive denominator it is enough to
show
\[
 Nk\sum_{r=k}^{N-1}\binom{N-1}{r}
 \le2^N(N-1)\binom{N-2}{k-1}.                              \tag{A.3b}\label{A.3b}
\]
For $k\le N/2$, use the full binomial sum and
$\binom{N-2}{k-1}\ge N-2$ away from the immediate endpoint $k=1$; for
$k\ge2,N\ge4$, the remaining comparison follows from
$Nk\le N^2/2\le2(N-1)(N-2)$.  The case $k=1$ is immediate.
For $k>N/2$, symmetry and monotonicity of the lower half of the binomial
row give
\[
 \sum_{r=k}^{N-1}\binom{N-1}{r}
 \le (N-k)\frac{N-1}{k}\binom{N-2}{k-1},
\]
and $N(N-k)\le2^N$ completes the proof.  The order $N=2$ is direct.

\subsection{Standard column-imbalance directions}

Let $\xi\in\mathbb R^{N+1}$ have coordinate sum zero and let $E(\xi)$ be
the standard embedding in \eqref{4.8}.  By $S_{N+1}$-invariance and
irreducibility of the standard module, it is enough to calculate the scalar
for the canonical vector $\xi=e_x-e_y$; the resulting ratio then holds for
every zero-sum $\xi$.  Distinguish the labels $x,y$.  The
signed two-label quotient has positive channels
\[
\begin{array}{c|l}
P_k&v=x,\ x,y\notin B,\quad 1\le k<N,\\
Q_k&v=x,\ y\in B,\quad 1\le k\le N,\\
R_k&v\notin\{x,y\},\ x\in B,\ y\notin B,\quad 1\le k<N,
\end{array}
\]
with the negative orientation obtained by exchanging $x,y$.  Use the channel
orders
\[
 \mathcal G=(P_1,\ldots,P_{N-1},R_1,\ldots,R_{N-1}),
 \qquad \mathcal Q=(Q_1,\ldots,Q_N).
\]
In these coordinates the operator is
\[
 H=\begin{pmatrix}S&C\\-D&Q\end{pmatrix},                 \tag{A.4s}\label{A.4s}
\]
where all four displayed blocks are nonnegative.  The nonzero signed rows
are
\[
 P_k\mapsto {\frac{k}{2N}}P_k+{\frac{N-k-1}{2N}}P_{k+1}
                   +{\frac{1}{2N}}Q_{k+1},                    \tag{A.5s}\label{A.5s}
\]
\[
\begin{aligned}
 Q_k\mapsto{}&{\frac{k^2-1}{2kN}}Q_k+{\frac{N-k}{2N}}Q_{k+1}
 -{\frac{k-1}{2kN}}P_{k-1}-{\frac{N-k}{2kN}}P_k\\
 &-{\frac{(k-1)^2}{2kN}}R_{k-1}
 -{\frac{(k-1)(N-k)}{2kN}}R_k,
\end{aligned}                                               \tag{A.6s}\label{A.6s}
\]
and
\[
\begin{aligned}
 R_k\mapsto{}&\left\{{\frac{k}{2N}}+{\frac{(k-1)(N-k)}{2kN}}\right\}R_k
 +{\frac{N-k-1}{2N}}R_{k+1}+{\frac{(k-1)^2}{2kN}}R_{k-1}\\
 &+{\frac{1}{2kN}}Q_k+{\frac{k-1}{2kN}}P_{k-1}
 +{\frac{N-k}{2kN}}P_k .
\end{aligned}                                               \tag{A.7s}\label{A.7s}
\]
Terms outside their physical ranges are absent.  Define the binomial row
\[
 \sigma(R_k)={\frac{\binom{N-2}{k-1}}{2^{N-2}}},\qquad
 \sigma(P_k)=0,\qquad \sigma(Q_k)=0,                       \tag{A.8s}\label{A.8s}
\]
and the physical reward
\[
\begin{aligned}
 \gamma(P_k)&=k d_k,\\
 \gamma(R_k)&=N d_k+{\frac{(N+1)(k-1)}{ k}}d_{k-1},\\
 \gamma(Q_k)&=-q_k,
 \quad q_k=(N-k)d_k+{\frac{(N+1)(k-1)}{ k}}d_{k-1}>0.
\end{aligned}                                               \tag{A.9s}\label{A.9s}
\]
The standard-sector scalar is
\[
 \Phi_N=\sigma(I-H)^{-1}\gamma.                            \tag{A.10s}\label{A.10s}
\]
This is the physical, rather than an auxiliary, reward.  Indeed the first
perturbation has feature coefficients
\[
 a_k={\frac{kd_k}{2(N+1)(N-1)}},\qquad
 b_k={\frac{Nd_k}{2(N+1)(N-1)}}
      +{\frac{(k-1)d_{k-1}}{2k(N-1)}},                         \tag{A.11s}\label{A.11s}
\]
whose $P,R,Q$ values are $a_k,b_k,a_k-b_k$.  The averaged second
perturbation annihilates the $a$ feature and has $b_k$ coefficient
$\sigma(R_k)/[2(N+1)]$.  Therefore
\[
 \boxed{{\frac{\Rtwo_n(E(\xi))}{\|\xi\|^2}}
 = {\frac{\Phi_N}{4(N+1)^2(N-1)}}.}                            \tag{A.12s}\label{A.12s}
\]

It remains to prove $\Phi_N>0$.  Put
\[
 R_S=(I-S)^{-1},\quad R_Q=(I-Q)^{-1},\quad
 A=R_SCR_QD,
\]
where the displayed rows give
$S\one\le(N-1)\one/N$ and
$(Q\one)_k=1/2-1/(2kN)$.  Thus both inverses are convergent Neumann series
and are entrywise nonnegative.  Also put
\[
 h=R_Qq,\qquad r_0=\gamma_S-Ch,\qquad f_0=R_Sr_0.          \tag{A.13s}\label{A.13s}
\]
Schur elimination of the bad channel gives the exact alternating-phase
identity
\[
 \Phi_N=\sigma(I+A)^{-1}f_0.                               \tag{A.14s}\label{A.14s}
\]

For $N\ge6$, define
\[
 W_1=2N^2d_1,\quad W_2=2Nd_1,\quad
 W_k={\frac{4N(k-1)}{ k}}d_{k-1}\quad(3\le k\le N).           \tag{A.15s}\label{A.15s}
\]
Then $(I-Q)W\ge q$.  At ranks one and two the cleared residuals are
$2(N^2-N+1)d_1$ and
\[
 {\frac{3N-4}{2}}d_1-{\frac{7(N-2)}{3}}d_2
 \ge {\frac{2(N-2)(N-6)}{3N}}.
\]
For $k\ge3$ the residual is $A_kd_{k-1}-B_kd_k$, with
\[
 A_k={\frac{(k-1)(3kN-2k^2-k+2)}{ k^2}},\qquad
 B_k={\frac{(3k+1)(N-k)}{ k+1}}.
\]
Using \eqref{A.3a}, it is at least
$2P(N,k)/[Nk^2(k+1)]$.  On setting $a=k-3$ and $m=N-k$,
\[
\begin{aligned}
P={}&a^4+a^3m+10a^3+5a^2m+37a^2+60a\\
   &+2am(m-1)+6m^2-22m+32>0.
\end{aligned}                                               \tag{A.16s}\label{A.16s}
\]
the last quadratic has discriminant $-284$.  Thus $h\le W$.  Direct
substitution now gives
\[
\begin{array}{c|c}
P_1&0\\
P_k&k(k-1)d_k/(k+1),\quad 2\le k<N\\
R_1&0\\
R_2&Nd_2+Nd_1/2\\
R_k&Nd_k+{\frac{k-1}{ k}}(N+1-2/k)d_{k-1},\quad 3\le k<N
\end{array}                                                \tag{A.17s}\label{A.17s}
\]
for the coordinates of $\gamma_S-CW$.  Hence $f_0\ge0$.

The bad exit and retention row sums satisfy
\[
 D\one={\frac{N-1}{2N}}\one,\qquad
 (Q\one)_k={\frac{1}{2}}-{\frac{1}{2kN}}.                         \tag{A.18s}\label{A.18s}
\]
Consequently $R_QD\one\le\one$.  With
\[
 z(P_k)=1/N,\qquad z(R_k)=2/(N+k),
\]
one has $(I-S)z\ge C\one$.  The $P_k$ gap is $1/(2N^2)$;
after $a=k-1,m=N-k$, the cleared $R_k$ gap is
\[
\begin{aligned}
Z={}&4a^4+14a^3m+18a^3+14a^2m^2+40a^2m+30a^2\\
 &+4am^3+24am^2+36am+20a+3m^3+8m^2+9m+4>0.
\end{aligned}                                               \tag{A.19s}\label{A.19s}
\]
Applying the positive inverses proves
\[
 A\one\le {\frac{2}{ N+1}}\one.                              \tag{A.20s}\label{A.20s}
\]

Two elementary barriers control the first phase.  The supersolution
$u(P_k)=4,u(R_k)=8N$ gives $0\le f_0\le8N\one$.  Indeed its
$P$ residual is $2(N+1)/N\ge kd_k$, while its $R$ residual exceeds the
upper bound $2(2N+1)/k$ for $\gamma(R_k)$ by
$2(4Nk-4N+1)/(Nk)$.  Separately, $q_k\le2N$ and
\eqref{A.18s} imply $h\le4N\one$.  For $N\ge7$, the subsolution
$\ell(P_k)=0,\ell(R_k)=2N$ then gives
\[
 f_0(R_k)\ge2N,\qquad \sigma f_0\ge2N.                    \tag{A.21s}\label{A.21s}
\]
At $k=1$ the margin is $2N-6-(N+1)\ge0$; at $k\ge2$ it is
$[3N-7-2k-4/N]/k\ge0$.

Set $c_N=2/(N+1)$.  Equations \eqref{A.20s}--\eqref{A.21s} and the
Neumann series for $(I+A)^{-1}$ yield, for $N\ge10$,
\[
 \Phi_N\ge2N-8N\sum_{j\ge1}c_N^j
 ={\frac{2N(N-9)}{ N-1}}>0.                                   \tag{A.22s}\label{A.22s}
\]
The remaining exact values are
\[
\begin{array}{c|c@{\qquad}c|c}
N&\Phi_N&N&\Phi_N\\ \hline
2&24/11&6&5758562957/248448224\\
3&261/40&7&141339691089527/4988552903680\\
4&343400/28657&8&15468663676289/466560376100\\
5&2268275/128288&9&19782952499295763/524622207176704.
\end{array}                                                \tag{A.23s}\label{A.23s}
\]
This proves the standard-sector sign in every order.  The supplementary
verifier independently reconstructs the quotient, normalization, barriers,
shifted polynomials, and every value in \eqref{A.23s}.

\subsection{Antisymmetric balanced directions}

Let $\delta=-\delta^{\mathsf T}$ have zero row sums and set
\[
 x(B,v)=\sum_{i\in B}\delta_{vi}.                            \tag{A.4}\label{A.4}
\]
Direct expansion of the active moves gives
\[
 \Delta h(B,v)=\frac{d_k}{2}x(B,v),\qquad k=|B|.             \tag{A.5}\label{A.5}
\]
The Poisson gradients satisfy
\[
 d_1>d_2>\cdots>d_{N-1}>0.                                  \tag{A.6}\label{A.6}
\]
To see this, couple adjacent complete rank chains by refreshing the same
bit.  Until their distinguished unequal bit is refreshed, the forcing
difference is $1/(X+1)-1/(X+2)>0$; adding a shared occupied bit decreases
that difference.  Summing the convergent Poisson series proves positivity
and strict monotonicity.

For any rank sequence $f_k$,
\[
 K_0\{f_kx(B,v)\}
 =\frac{kf_k+(N-k-1)f_{k+1}}{2N}x(B,v).                      \tag{A.7}\label{A.7}
\]
Thus $G\Delta h=r_kx$, where $r_N=0$ and
\[
 (2N-k)r_k-(N-k-1)r_{k+1}=Nd_k.                             \tag{A.8}\label{A.8}
\]
Backward induction using \eqref{A.6} gives
\[
 0\le r_{k+1}<r_k\le\frac{N}{N+1}d_k.                       \tag{A.9}\label{A.9}
\]

Put $T=\|\delta\|_F^2$.  Sampling without replacement and applying
$\Delta$ a second time yields
\[
\begin{aligned}
 \E_k[\Delta(r_kx)]={\frac{T}{2nN}}\Big[&
 {\frac{k(N-k)}{ N-1}}(r_k-r_{k+1})+(N-k)r_{k+1}\\
 &+{\frac{(k-1)(N-k+1)}{ N-1}}(r_{k-1}-r_k)
 +(N-k+1)r_k\Big].
\end{aligned}                                               \tag{A.10}\label{A.10}
\]
The coefficient of the absent $r_0-r_1$ term is zero.  Every displayed term
is nonnegative by \eqref{A.9}, and the total is strict when $T>0$.
Averaging with \eqref{A.1} proves positivity of the antisymmetric sector for
every $N\ge2$.

\subsection{Symmetric balanced directions}

For a symmetric balanced $\delta$, put
\[
 x(B,v)=\sum_{i\in B}\delta_{vi},\qquad
 z(B)=\sum_{w\ne i\in B}\delta_{wi}.                       \tag{A.11}\label{A.11}
\]
Functions in this sector have the form $a_kx+b_kz$.  Transpose the exact
two-channel rank operator and split its positive and negative channels:
\[
 H=K^{\mathsf T}=\begin{pmatrix}S&C\\-D&Q\end{pmatrix}.     \tag{A.12}\label{A.12}
\]
Every displayed block is entrywise nonnegative.  The nonzero entries are
\[
\begin{aligned}
 S_{k,k}&={\frac{k}{2N}},&S_{k,k-1}&={\frac{N-k}{2N}},\\
 Q_{k,k}&={\frac{N(k-2)+k}{2kN}},
 &Q_{k,k-1}&={\frac{N-k-1}{2N}},\\
 Q_{k,k+1}&={\frac{k(k-1)}{2(k+1)N}},
 &D_{k,k-1}&={\frac{1}{ N}},
\end{aligned}                                               \tag{A.13}\label{A.13}
\]
and
\[
 C_{k-1,k}={\frac{k-1}{2kN}},\qquad C_{k,k}={\frac{N-k}{2kN}}.      \tag{A.14}\label{A.14}
\]
The $a$ channel is indexed by $1\le k<N$ and the $b$ channel by
$2\le k<N$; entries outside these ranges are zero.  The positive source is
\[
 s_k^a={\frac{d_k}{2}},\qquad s_k^b={\frac{d_{k-1}}{2k}},            \tag{A.15}\label{A.15}
\]
and the good/bad reward is
\[
 g_k^a={\frac{\binom{N-2}{k-1}}{2^{N-1}(N+1)}},\qquad
 g_k^b=-{\frac{\binom{N-3}{k-2}}{2^{N-2}(N+1)}}.                \tag{A.16}\label{A.16}
\]

The normalization connecting this quotient to the physical Hessian is
exact, not merely sign preserving.  If $(a,b)=(I-K)^{-1}s$, set
$b_1=a_N=b_N=0$ and
\[
 r_k=a_k-\frac{2(k-1)}{N-2}b_k,
\]
then the fixed-count orbit identities give
\[
 \frac{\Rtwo_n(\delta)}{\|\delta\|_F^2}
 =\sum_{k=1}^{N-1}\pi_k\frac{N-k}{(N-1)(N+1)}r_k
 =g^{\mathsf T}(I-K)^{-1}s=:\mathcal S_N.                  \tag{A.17}\label{A.17}
\]
Here is the all-order labelled normalization.  Let
$c=1/[nN2^{N-1}]$ and write $\mu=\nu_0\Delta$.  For an output state $(B,v)$
of rank $k$, the four predecessor families contribute, respectively,
\[
 {\frac{c}{2}}kx(B,v),\quad {\frac{c}{2}}(k-1)x(B,v),\quad
 {\frac{c}{2}}(N-k)x(B,v),\quad {\frac{c}{2}}(N-k+1)x(B,v).
\]
Thus their sum is the exact labelled current
$\mu(B,v)=x(B,v)/(n2^{N-1})$, including the boundary ranks.  Define
$f(B,v)=a_{|B|}x(B,v)+b_{|B|}z(B)$.  Conditional on a uniform rank-$k$
orbit, sampling without replacement gives, with
$R_v=\sum_i\delta_{vi}^2$,
\[
 \E[x^2\mid v]={\frac{k(N-k)}{ N(N-1)}}R_v,\qquad
 \E[xz\mid v]=-{\frac{2k(k-1)(N-k)}{ N(N-1)(N-2)}}R_v.       \tag{A.17a}\label{A.17a}
\]
For the cross term, separate the two coincidences in the ordered triple
from the all-distinct case; the row and column sums reduce the resulting
full sums to $-2R_v$.  Summing $R_v$ over $v$ gives
$\|\delta\|_F^2$.  Therefore
\[
 {\frac{\nu_0\Delta f}{\|\delta\|_F^2}}
 =\sum_{k=1}^{N-1}\pi_k{\frac{N-k}{(N-1)(N+1)}}
 \left(a_k-{\frac{2(k-1)}{ N-2}}b_k\right),                  \tag{A.17b}\label{A.17b}
\]
which is the first expression in \eqref{A.17}.  The identities
\[
 \pi_k{\frac{N-k}{ N-1}}
 ={\frac{\binom{N-2}{k-1}}{2^{N-1}}},\qquad
 {\frac{k-1}{ N-2}}\binom{N-2}{k-1}=\binom{N-3}{k-2}
\]
show that the coefficients in \eqref{A.17b} are exactly $g_k^a,g_k^b$.
Hence the orbit average is
$g^{\mathsf T}(a,b)=g^{\mathsf T}(I-K)^{-1}s$.  This is an all-order labelled
derivation; the exact replay independently reconstructs it for
$4\le n\le12$ before taking a sign.

Put $R_S=(I-S)^{-1}$ and $R_Q=(I-Q)^{-1}$.  Here $S$ has row sums at most
$1/2$, while the displayed tridiagonal $Q$ is substochastic and every
communicating class reaches a strict boundary row.  Hence both inverses are
convergent entrywise-nonnegative Neumann series.  Define
\[
\begin{aligned}
 W&=R_Q(-g^b),& r&=g^a-CW,& f_0&=R_Sr,\\
 A&=R_SCR_QD,& \ell&=(s^a)^{\mathsf T}-(s^b)^{\mathsf T}R_QD,&
 \mathcal D&=(s^b)^{\mathsf T}W.
\end{aligned}                                               \tag{A.18}\label{A.18}
\]
Solving the bad equation first gives the exact Schur identity
\[
 \mathcal S_N=\ell(I+A)^{-1}f_0-\mathcal D.                 \tag{A.19}\label{A.19}
\]
One factor of the nonnegative matrix $A$ is one completed visit to the bad
channel.  The inverse in \eqref{A.19} therefore alternates successive
re-entries, so a first-excursion estimate alone would not suffice.

We now give the full phase certificate.  Put
\[
 v=R_Sg^a,\qquad v_k=t_k g_k^a,\qquad
 t_1={\frac{2N}{2N-1}},\quad
 t_k={\frac{2N+(k-1)t_{k-1}}{2N-k}}.                           \tag{A.20}\label{A.20}
\]
Induction gives
\[
 t_k\le {\frac{N^2-k}{(N-2)(N-k)}}.                            \tag{A.21}\label{A.21}
\]
The base difference is $(5N-1)/[(N-2)(2N-1)]$; after one step and the
substitution $m=N-k$, the remaining numerator is
$N^3-N^2+2Nm^2+2Nm+2m^3+3m^2+m>0$.

For $N\ge24$, let $\overline W=(7N/25)(-g^b)$.  Direct substitution shows
\[
 (I-Q)\overline W\ge-g^b,\qquad
 C\overline W\le {\frac{14}{25}}g^a.                           \tag{A.22}\label{A.22}
\]
For the first inequality, after clearing positive binomial factors the
interior residual is
\[
 P_N(k)=21N^2k+14N^2-64Nk^2-57Nk+50k^3+36k^2-14k.          \tag{A.23}\label{A.23}
\]
For $N\ge25$, $\operatorname{disc}_kP_N=-8B_N$, where the coefficients of
$B_{25+M}$ are
\[
 (5733,736211,37934833,982793213,12893444604,
 70866894144,41021531998).                                 \tag{A.24}\label{A.24}
\]
Thus $P_N$ has one real, negative root and is positive for $k\ge0$; at
$N=24$ its exact minimum over the physical ranks is $24$.  Applying $R_Q$
in \eqref{A.22} yields
\[
 {\frac{11}{25}}v\le f_0\le v.                                \tag{A.25}\label{A.25}
\]

For $2\le k<N$, define
\[
 \widehat h_k={\frac{k}{ N-2}}g^a_{k-1}.                       \tag{A.26}\label{A.26}
\]
Equations \eqref{A.13} and \eqref{A.21} give
$(I-Q)\widehat h\ge Dv$ and
\[
 C\widehat h\le c_Ng^a,\qquad
 c_N={\frac{2N-5}{2N(N-2)}}.
\]
Indeed, the interior ratio in the first comparison is
$(N^2-k+1)/[N(N-2)(N-k+1)]$, and the largest row ratio in the second is
at rank $N-2$.  At $k=2$ the first boundary residual is
$(N+1)/[N(N-2)]$, equal to the interior formula; at $k=N-1$ it is
$(N^2+1)/[2N(N-2)]$, which exceeds the interior value
$(N^2-N+2)/[2N(N-2)]$.  Thus \eqref{A.21} dominates every boundary row as
well.  Consequently
\[
 Av\le c_Nv.                                               \tag{A.27}\label{A.27}
\]

For the left occupation debt, put
\[
 Y_k={\frac{2(k+1)}{3k(k-1)}}\qquad(2\le k<N).                 \tag{A.28}\label{A.28}
\]
With $B=Q^{\mathsf T}$, an interior substitution gives
\[
 [(I-B)Y]_k-{\frac{1}{ k(k-1)}}
 ={\frac{4Nk+2N+2k^3-k^2-5k}{3Nk^2(k-1)(k+1)}}>0,             \tag{A.29}\label{A.29}
\]
while the boundary residuals at $k=2$ and $k=N-1$ are, respectively,
$(5N+7)/(18N)>0$ and
$(N^2-N+2)/[3N(N-2)(N-1)^2]>0$.  Since
$s^b_k\le1/[k(k-1)]$, this proves
$(s^b)^{\mathsf T}R_Q\le Y^{\mathsf T}$ and hence
$\ell\ge\overline\ell>0$, where
\[
 \overline\ell_j={\frac{3Nj+3N-8j-10}{3Nj(j+1)}}.             \tag{A.30}\label{A.30}
\]
The absent bad state above the top rank only increases the final margin.

The bad debt and first phase now satisfy
\[
 \mathcal D\le\sum_{j=1}^{N-2}
 {\frac{4(j+2)}{3(j+1)(N-2)}}g_j^a,\qquad
 \ell f_0\ge{\frac{11}{25}}\sum_{j=1}^{N-1}
 \overline\ell_jt_jg_j^a.                                \tag{A.31}\label{A.31}
\]
Define
\[
\begin{aligned}
 \beta_N&=\max_{1\le j\le N-2}
 {\frac{4(j+2)}{3(j+1)(N-2)}}
 \left({\frac{11}{25}}\overline\ell_jt_j\right)^{-1},\\
 \varepsilon_N&={\frac{25}{11}}{\frac{c_N}{1-c_N}}.
\end{aligned}                                              \tag{A.32}\label{A.32}
\]
Exact rational recurrence \eqref{A.20} gives
$\beta_N+\varepsilon_N<1$ for $40\le N\le287$; its smallest
margin is
\[
1-\beta_{40}-\varepsilon_{40}
={\frac{639304267467075678841}{115369588296792467144716}}>0.  \tag{A.33}\label{A.33}
\]
This is a finite exact calculation.

For $N\ge288$, induction in \eqref{A.20} gives
$t_j\ge N/(N-j+\sqrt{N/2})$.  Indeed, substituting the inductive bound leaves
the numerator
$2u(N-j+u)+u-N$, where $u=\sqrt{N/2}$; it is minimized at $j=N-1$ and
equals $3u>0$.  Since $u\le N/24$ in this range,
$t_j\ge24N/(25N-24j)$.  The bound
$\beta_N\le19/20$ is then equivalent, after clearing positive factors, to
$G_N(j)>0$, where
\[
\begin{aligned}
G_N(j)={}&1881N^2j+1881N^2-6250Nj^2-21278Nj-10032N\\
         &+6000j^3+12000j^2+10032j+12540.
\end{aligned}                                               \tag{A.34}\label{A.34}
\]
Its discriminant is $-500D_N$, and the coefficients of $D_{288+M}$ are
\[
\begin{gathered}
43034652243,71940715263252,50075984929826764,\\
18577025353519361184,3873653773133773223808,\\
430447522448513182675968,
19913301794000751100222464.
\end{gathered}                                             \tag{A.35}\label{A.35}
\]
Thus $G_N$ has one negative real root and is positive on every physical
rank.  Also $\varepsilon_N\le1/20$ for $N\ge46$.  Therefore
$\beta_N+\varepsilon_N<1$ for all $N\ge288$.

Finally, \eqref{A.25}, \eqref{A.27}, and positivity of $\ell$ give
\[
 \left|\ell(I+A)^{-1}f_0-\ell f_0\right|
 \le\varepsilon_N\ell f_0.
\]
Together with \eqref{A.31} this yields
\[
 \mathcal S_N\ge(1-\beta_N-\varepsilon_N)\ell f_0>0       \tag{A.36}\label{A.36}
\]
for every $N\ge40$.  The finitely many orders $3\le N\le39$ are solved
exactly over $\mathbb Q$; the first values are
$\mathcal S_3=3/208$ and $\mathcal S_4=359/26660$.

The supplementary verifier reconstructs every rational boundary solve,
cleared polynomial, shifted coefficient list, and the labelled
normalizations \eqref{A.12s} and \eqref{A.17}.  Combining the three
subsections proves Theorem~\ref{thm:sectors}.
